% docs/physics/latex/physics_reference.tex - the bound book "SIMC++ Physics Reference".
% Single source: the chapters ARE the note files of this directory (\simcinputnote inputs a complete note without its
% preamble). Part I is the common note 00, Part II the reaction notes in the order of notes.json
% (physics_reference_notes.tex is written by build_physics_pdf.py). Do not put physics text here.
% Build: cmake --build <build> --target physics_pdf (see docs/physics/README.md).
\documentclass[book]{simcnote}
\simcsetmeta{SIMC++ Physics Reference}%
  {The physics of the SIMC++ reactions: generation, kinematics, cross sections, radiation and the plug-ins that
   implement them}%
  {\simcbuildversion}{\simcbuilddate}{SIMC++ documentation}
\begin{document}
\simcphysfrontmatter
\chapter*{About this book}\addcontentsline{toc}{chapter}{About this book}\markboth{About this book}{}
\simcabstract{One chapter per reaction of SIMC++, preceded by the elements they share.}
Each chapter is a \emph{physics note}: it describes one reaction as SIMC++ builds it (the reaction and physics-model
plug-ins, how the pipeline selects them, the stages they use), then every formula the code evaluates, the inputs and
outputs of the reaction, its validation and its Fortran origin. Each displayed equation names the C++ function that
computes it and, second, the line of the Fortran SIMC (upstream commit \texttt{59ccdf6}) it was translated from.

The same files are also built as standalone PDFs, one per note, in \texttt{docs/physics/pdf/}; a copy of each
reaction note sits in the source folders of its plug-ins (\texttt{src/physics/<reaction>/}), next to the code it
describes. Chapter~\ref{N00-note} (Part~I) defines the notation, the glossary and the common physics used by all the
other chapters. The user manual (chapter 5, Reactions) explains how to run each reaction; the developer manual
(chapter 4) explains how to add a reaction or a physics model, which comes with a new chapter of this book written
from \texttt{docs/physics/latex/template\_reaction.tex}.
\simcmainmatter
\simcpart{Common elements}{The plug-in architecture, the event chain, notation and glossary, and the physics shared by
all reactions: kinematics, weighting and normalisation, target effects, Coulomb correction and radiation.}
\simcinputnote{00_common}
\simcpart{Reactions}{One chapter per reaction: the SIMC++ components, kinematics and generation, the cross-section
model, radiation, inputs and outputs, validation and the Fortran origin.}
\input{physics_reference_notes}
\end{document}
